% =====================================================================
%  KẾT LUẬN
% =====================================================================
\chapter*{Kết Luận}
\addcontentsline{toc}{chapter}{Kết Luận}
\markboth{Kết Luận}{}
\label{ch:ket-luan}

\section*{Kết quả đạt được}

Tiểu luận đã đi trọn một vòng từ lịch sử, lý thuyết, cài đặt, thực nghiệm đến triển khai cho ba họ mô hình trên bảy tập dữ
liệu thật:
\begin{itemize}
  \item \textbf{Lịch sử AI} được hệ thống hóa qua mười sáu mốc, cho thấy mỗi mô hình hiện đại là lời đáp cho một thất bại cụ
    thể trước đó: lan truyền ngược cho giới hạn XOR, học máy thống kê cho tính giòn của hệ chuyên gia, LSTM cho triệt tiêu
    đạo hàm, ResNet cho thoái hóa mạng sâu.
  \item \textbf{Học máy cơ bản và MLP:} trên dữ liệu y tế mất cân bằng, MLP đạt F1 lớp dương \SL{ch2_mlp_torch/diabetes/f1},
    cao nhất trong bảy mô hình; trên giá nhà Việt Nam, Random Forest vẫn nhỉnh nhất ($R^2 = \SL{ch2_sklearn/house/rf/r2}$);
    trên văn bản, MLP chính xác nhất nhưng quá khớp nhanh. MLP cài đặt bằng NumPy thuần đạt kết quả ngang Keras và PyTorch.
  \item \textbf{CNN:} ResNet-18 huấn luyện từ đầu đạt \SL{ch3_svhn/resnet18/accPct}\,\% trên SVHN và
    \SL{ch3_eurosat/resnet18/accPct}\,\% trên EuroSAT; MobileNetV2 đạt độ chính xác sát nút với số tham số ít hơn khoảng
    5 lần, nhưng không nhanh hơn trên GPU ở lô 1. Mạng CNN chạy bằng các hàm NumPy viết tay cho cùng dự đoán với PyTorch.
  \item \textbf{RNN:} trên KKBox, các tế bào có cổng vượt Simple RNN; trên giá FPT, cả bốn kiến trúc chỉ ngang mốc “giữ nguyên
    giá” dù $R^2$ trên giá đạt \SL{ch4_torch/stock/lstm/r2}. LSTM viết tay bằng NumPy tái tạo đúng đầu ra của PyTorch.
  \item \textbf{Triển khai:} một Web App Flask sáu phân hệ phục vụ mọi mô hình trên CPU với thời gian phản hồi vài đến vài
    chục mili-giây, có bản đồ CAM giải thích dự đoán ảnh.
\end{itemize}

\section*{Bài học rút ra}

\begin{enumerate}
  \item \textbf{Chọn mô hình theo dạng dữ liệu, không theo độ “hiện đại”.} Random Forest vẫn là đối thủ khó vượt trên dữ liệu
    bảng; CNN thắng trên ảnh vì cấu trúc không gian; RNN chỉ có lợi khi chuỗi thật sự chứa tín hiệu theo thời gian.
  \item \textbf{Một con số đẹp cần một mốc so sánh.} $R^2 = 0{,}99$ trên giá cổ phiếu hay Accuracy 91,5\,\% trên dữ liệu tiểu
    đường đều đạt được bởi mô hình ngây thơ. Mọi đánh giá trong báo cáo vì vậy luôn đi kèm mốc tham chiếu và chỉ số phù hợp
    với bài toán (F1 lớp hiếm, $R^2$ trên lợi suất).
  \item \textbf{Tiền xử lý quyết định nhiều hơn kiến trúc.} Chuẩn hóa theo mức giá khiến mọi RNN thất bại khi giá vượt vùng
    huấn luyện; chuẩn hóa theo từng cửa sổ sửa được ngay. Ngưỡng phân loại thay đổi F1 nhiều hơn việc đổi mô hình.
  \item \textbf{FLOPs không phải độ trễ.} MobileNetV2 rẻ hơn nhiều lần về phép tính nhưng chậm hơn ResNet-18 trên GPU ở lô 1;
    hiệu năng thực phụ thuộc phần cứng đích.
  \item \textbf{Viết tay để hiểu, dùng thư viện để chạy.} Các bản NumPy giúp kiểm chứng từng công thức; PyTorch trên GPU nhanh
    hơn nhiều bậc khi huấn luyện thật.
\end{enumerate}

\section*{Hướng phát triển}

Các hướng tiếp theo bám sát lộ trình của học phần: thay RNN bằng Transformer cho chuỗi dài; bổ sung dữ liệu ngoài chuỗi giá
(tin tức, báo cáo tài chính) và thông tin giao dịch thuê bao cho bài toán rời bỏ; dùng học chuyển giao với trọng số tiền huấn
luyện cho ảnh; và kết nối các mô hình trong báo cáo với đồ thị tri thức, RAG và hệ đa tác tử ở các chương sau của học phần.
